\section{Related work}\label{sec:related}

ClinGen maps predictor scores to ACMG/AMP evidence strengths through local likelihood ratios estimated against ClinVar
labels with a one-sided bootstrap bound \cite{Pejaver2022}, extended to more predictors \cite{Bergquist2025}, to assay
scores \cite{Brnich2020,Fayer2021,Zeiberg2025} and to gene- and domain-specific calibration
\cite{Tejura2024,Chen2026,Badonyi2025}, the last with sign-test bounds on resampled likelihood ratios \cite{Badonyi2025}.
Split conformal prediction gives finite-sample coverage \cite{Vovk2005}; its class-conditional (Mondrian) form controls
coverage within each class \cite{Vovk2005,Sadinle2019,Ding2023}; order-statistic thresholds give training-conditional
control \cite{Vovk2013,Tong2018,Park2020}; keeping one member per group handles clustered data \cite{Dunn2023}; and
abstaining classifiers trade coverage against error \cite{Chow1970}. Selecting a threshold with exact binomial bounds and
a union bound underlies selective classification with guaranteed risk \cite{Geifman2017} and risk-controlling prediction
sets \cite{Bates2021rcps}; fixed-sequence testing over an ordered grid is the Learn-then-Test template
\cite{Angelopoulos2025ltt}. Harrizi et al.\ reported split-conformal prediction sets for missense predictors on ClinVar
at 90\% coverage \cite{Harrizi2026}; \textsc{CertiVar} instead controls coverage within each class, calibrates on assay
labels and certifies the likelihood ratio of the calls. Large-sample intervals for a ratio of binomial proportions
\cite{Katz1978,Koopman1984} carry no finite-sample guarantee (\cref{sec:evidence} compares Katz intervals). Label-shift methods
adjust posteriors or conformal weights to an estimated target prior \cite{Saerens2002,Lipton2018,Podkopaev2021};
class-conditional quantities need none, provided class-conditional invariance holds, which \cref{sec:method} tests.
